% Chapter 5, typeset from lectures/L05: the README is the opener and the closing review, the
% lecture's Appendix A is the chapter's one section of material, and Appendix B its exercises.

\chapter{Timers}
\label{c5:ch}

{\large\itshape Letting the hardware keep time, so the CPU is free to do something else.\par}

\begin{chapterbox}{In this chapter}
\begin{itemize}
  \item Why \code{_delay_ms()} doesn't scale, and what a hardware timer gives you instead.
  \item The prescaler math: from CPU clock to tick period to interrupt count.
  \item Normal Mode vs.\ CTC Mode.
  \item Worked example: non-blocking LED blink using Timer 0 (Normal Mode).
  \item Timer 1 in CTC Mode: hitting an arbitrary period directly.
\end{itemize}
\end{chapterbox}

Every delay so far has been a busy-wait. This chapter explains why a busy-wait delay
(\code{_delay_ms}) is unsuitable once a program needs to do more than one thing at a time, and
replaces it with the chip's hardware timers. You will compute a timer's tick period from the CPU
clock and a chosen prescaler, work out how many overflow interrupts an 8-bit timer needs for a
given delay, and learn the practical difference between Normal Mode and CTC Mode, and when you'd
reach for each. Along the way the chapter returns to \code{volatile}, and to why a variable shared
between an ISR and \code{main()} needs it.

Read \secref{c5:sec:timers} first: it is the chapter's core material, with its worked examples.
Then work the exercises in \secref{c5:sec:exercises}. The Reflection and Calculation exercises,
and the questions some Code exercises ask, are answered in Appendix~\ref{app:answers}; the Code
exercises themselves have no printed solution, because the board tells you whether a program does
what the exercise asks.

% Section 5.1, from lectures/L05/appendix/a_timers.md. The lecture's two circuit images are copies
% of Chapter 4's potentiometer circuit, so neither is printed again: both circuits are described
% against Figure 4.1 instead.

\section{Timers}
\label{c5:sec:timers}\label{c5:app:a}

\subsection{Why timers}
\label{c5:a:1}
In the previous examples, a function \code{delay_ms()} has been used to generate delays between LED
blinks and PWM pulses:

\begin{ccode}
/**
 * @brief Generate delay.
 *
 * @param[in] ms Delay duration in ms.
 */
static void delay_ms(const uint16_t ms)
{
    for (uint16_t i = 0U; i < ms; ++i)
    {
        _delay_ms(TICK_1MS);
    }
}
\end{ccode}

\noindent This function works as intended, but it has some flaws:
\begin{itemize}
  \item \code{_delay_ms()} blocks: the CPU sits in a busy-loop doing nothing else for the duration.
    Fine for a two-line demo, not for anything that also needs to poll a button, read a sensor, or
    respond to interrupts while ``waiting.''
  \item To avoid CPU blocks, a timer circuit can be used that flags when the given duration has
    passed, i.e.\ when timeout occurs.
  \item A timer is hardware that counts up on its own, driven by the CPU clock (divided down by a
    \textbf{prescaler}), completely independent of your code. Configure it once and it can either
    be polled, or, more usefully, fire an interrupt when it reaches a target value, letting you
    build non-blocking delays and precise periodic events.
  \item The ATmega328P has three: \textbf{Timer 0} and \textbf{Timer 2} are 8-bit (count
    0--255), \textbf{Timer 1} is 16-bit (count 0--65535). Their registers follow a consistent
    naming pattern; swap the digit (\code{0}/\code{1}/\code{2}) and the names carry over,
    although, as \secref{c5:a:2} and \secref{c5:a:5} show, not every bit pattern does (Timer 2's
    prescaler bits, and the bit that selects CTC mode on Timer 1).
\end{itemize}

\subsection{The math: prescaler \texorpdfstring{$\rightarrow$}{->} tick time
  \texorpdfstring{$\rightarrow$}{->} interrupt count}
\label{c5:a:2}
Pick a \textbf{prescaler} (a clock divider) to get a manageable tick rate. At 16\,MHz with a
\textdiv1024 prescaler:

\begin{console}
tick frequency = 16,000,000 / 1024 = 15,625 Hz
tick period    = 1 / 15,625 = 0.064 ms
\end{console}

\noindent For an 8-bit timer, one full count (0$\rightarrow$255, i.e.\ 256 ticks) triggers one
\textbf{overflow}. To get a specific delay, work out how many overflows you need:

\begin{console}
ticks needed     = desired_delay_ms / tick_period_ms
overflows needed = ticks needed / 256          (8-bit timer)
\end{console}

\noindent E.g.\ a 100\,ms delay at 0.064\,ms/tick needs $\sim$1562 ticks $\approx$ 6 overflows. You
don't have to hit this exactly with a single hardware compare; instead count overflow interrupts
in software and act once you've counted enough. That's the pattern in the example below.

Prescaler bit combinations for Timer 0 / Timer 1 are shown below:

\begin{narrowtable}{@{}lrr@{}}
  \thead{Prescaler bits} & \thead{Divisor} & \thead{Tick @ 16\,MHz} \\
  \midrule
  \code{001} & 1 & 62.5\,ns \\
  \code{010} & 8 & 0.5\,\textmu s \\
  \code{011} & 64 & 4\,\textmu s \\
  \code{100} & 256 & 16\,\textmu s \\
  \code{101} & 1024 & 64\,\textmu s \\
\end{narrowtable}

\noindent Timer 2's prescaler bit values differ somewhat:

\begin{narrowtable}{@{}lrr@{}}
  \thead{Prescaler bits} & \thead{Divisor} & \thead{Tick @ 16\,MHz} \\
  \midrule
  \code{001} & 1 & 62.5\,ns \\
  \code{010} & 8 & 0.5\,\textmu s \\
  \code{011} & 32 & 2\,\textmu s \\
  \code{100} & 64 & 4\,\textmu s \\
  \code{101} & 128 & 8\,\textmu s \\
  \code{110} & 256 & 16\,\textmu s \\
  \code{111} & 1024 & 64\,\textmu s \\
\end{narrowtable}

\subsection{Normal Mode vs.\ CTC Mode}
\label{c5:a:3}
There are two modes that will be covered in this chapter:
\begin{itemize}
  \item \textbf{Normal Mode}: the timer always counts to its full range (255 or 65535) before
    overflowing and firing \code{TIMERn_OVF_vect}.
  \item \textbf{CTC Mode} (Clear Timer on Compare): the timer counts up to a value \emph{you}
    choose (written to \code{OCRnA}), fires \code{TIMERn_COMPA_vect}, and resets to 0, useful when
    the period you want doesn't cleanly divide into 256 or 65536 ticks.
\end{itemize}
Both need the same steps: set the prescaler (which also starts the timer), enable the timer's own
interrupt in \code{TIMSKn}, and enable interrupts globally. CTC additionally needs the compare
value and the mode bit.

\subsection{Non-blocking blink using Timer 0 (Normal Mode)}
\label{c5:a:4}
LED on D9 toggles approximately every 500\,ms (31 overflows of 16.384\,ms each, 508\,ms); no
\code{_delay_ms()} anywhere, so the rest of \code{main()}'s loop stays free to do other work. The
circuit is Chapter~\ref{c2:ch}'s first one, Figure~\ref{c2:fig:ledbutton}; its button on D13 isn't
used here.

\begin{ccode}
#include <stdint.h>

#include <avr/interrupt.h>
#include <avr/io.h>

/** GPIO pins. */
#define LED1 1U // D9 -> PORTB1

/** GPIO operations. */
#define LED1_TOGGLE PINB = (1U << LED1) // Toggle LED1.

/** Time parameters. */
#define TICK_PERIOD_MS 0.064F // Tick period, prescaler 1024.
#define TICK_MAX 256U         // Ticks per overflow when using an 8-bit timer.
#define TIMEOUT_MS 500U       // Timeout in ms.

/** Limit parameters; ticks and overflows needed for the desired timeout. */
#define OVF_PERIOD_MS (TICK_PERIOD_MS * TICK_MAX)            // Time between each overflow in ms.
#define OVF_MAX (uint8_t)(TIMEOUT_MS / OVF_PERIOD_MS + 0.5F) // Timeout measured in overflows.

/**
 * @brief Set up system.
 */
static void setup(void)
{
    // Configure LED1 as output.
    DDRB = (1U << LED1);

    // Enable overflow interrupt for timer 0, use prescaler 1024.
    TCCR0B = (1U << CS00) | (1U << CS02);
    TIMSK0 = (1U << TOIE0);
    sei();
}

/**
 * @brief Count overflows, toggle LED1 on timeout.
 */
ISR(TIMER0_OVF_vect)
{
    static volatile uint8_t ovf_count = 0U;

    // Increment the overflow counter, toggle LED1 and clear the counter on timeout.
    if (OVF_MAX <= ++ovf_count)
    {
        LED1_TOGGLE;
        ovf_count = 0U;
    }
}

/**
 * @brief Application entry point.
 *
 * @return 0 on termination of the program (should never occur).
 */
int main(void)
{
    setup();
    while (1) {}
    return 0;
}
\end{ccode}

\code{OVF_PERIOD_MS} and \code{OVF_MAX} are the calculation macros covered in \secref{c1:a:12}:
\code{OVF_PERIOD_MS} parenthesizes its whole body the way \code{TICKS_NEEDED} does there, and
\code{OVF_MAX} gets the same safety from its leading \code{(uint8_t)} cast instead: the cast
applies to the parenthesized division that follows it, and a cast binds tighter than any
arithmetic or comparison operator, so the expanded macro still acts as a single operand wherever
it's used.

These macros use \code{float}, but they cost nothing at runtime: every operand is a constant, so the
compiler computes the result while compiling and only the final integer ends up in the program.
\secref{c4:a:5}'s advice to avoid \code{float} applies to math on values only known at runtime,
like an ADC reading.

\code{ovf_count} here is a function-local \code{static} variable, scoped entirely inside the ISR:
nothing outside \code{ISR(TIMER0_OVF_vect)} ever touches it. That means \code{volatile} isn't
actually load-bearing on this specific variable: the compiler can already see every read and write
to \code{ovf_count} within this one function body, and nothing else can change it between two of
them. Leaving \code{volatile} on here is a harmless habit. The moment to actually rely on it is the
moment code outside this ISR needs to read \code{ovf_count} too, above all a loop in \code{main()}
that waits for it to change, the situation the \code{blink_state} of \secref{c3:a:4} was in.

\subsection{Timer 1 in CTC Mode (16-bit, arbitrary period)}
\label{c5:a:5}
Same idea, but Timer 1's extra range plus CTC mode let you hit a period directly, without stacking
overflow counts except for very long delays. Register names swap the \code{0} for a \code{1}; CTC
needs one extra bit, \code{WGM12}, to select CTC instead of Normal Mode.

The counter, \code{TCNT1}, counts from \code{0} up to the value you write to \code{OCR1A}, so it
takes \code{(OCR1A + 1)} ticks to reach it and fire a compare-match interrupt. To hit a specific
timeout, work out how many ticks that timeout needs, then subtract one:

\begin{console}
ticks needed = timeout_ms / TICK_PERIOD_MS
OCR1A        = ticks needed - 1
\end{console}

\noindent E.g.\ for a 100\,ms timeout at a 0.064\,ms tick (\textdiv1024 prescaler):
\code{100 / 0.064 = 1562.5} ticks needed, exactly halfway between 1562 ticks (99.968\,ms) and 1563
(100.032\,ms), which are equally close. Rounding up gives 1563, so \code{OCR1A = 1563 - 1 = 1562}.
The macros below make the same calculation in floating point and land on \code{1561} instead, one
tick short: \code{0.064F} can't be stored exactly and is a hair above \code{0.064}, so the
division gives a hair under \code{1562.5}, and the rounding goes down. On a half tick either answer
is as good as the other; a timeout that doesn't land on one, like the two seconds of
Exercise~\ref{c5:ex:6}, comes out the same either way:

\begin{ccode}
#define TIMEOUT_MS 100U                            // Timeout in ms.
#define TICK_PERIOD_MS 0.064F                      // Tick period, prescaler 1024.
#define CMP_MAX (TIMEOUT_MS / TICK_PERIOD_MS - 1U) // Compare value.
#define TICK_MAX (uint16_t)(CMP_MAX + 0.5F)        // Rounded compare value.

TCCR1B = (1U << CS10) | (1U << CS12) | (1U << WGM12); // Set 1024 prescaler, use CTC mode.
OCR1A  = TICK_MAX;                                    // Set timeout to 100 ms.
TIMSK1 = (1U << OCIE1A);                              // Enable compare-match interrupt.
\end{ccode}

\noindent Timer 1's 16-bit range means \code{1561} fits directly; Timer 0's 8-bit range (max
\code{255}) couldn't hold a value that large, which is exactly why \secref{c5:a:4} had to count
multiple overflows in software instead.

Two LEDs on D9/D10 alternate once every 100\,ms while enabled; a button on D13 toggles them on and
off entirely, cutting the LEDs immediately rather than waiting for the next compare match. The
circuit is the one from \secref{c2:a:6}, Figure~\ref{c2:fig:twoleds}.

\begin{ccode}
#include <stdbool.h>
#include <stdint.h>

#include <avr/interrupt.h>
#include <avr/io.h>

/** GPIO pins. */
#define LED1 1U // D9  -> PORTB1.
#define LED2 2U // D10 -> PORTB2.
#define BTN1 5U // D13 -> PORTB5.

/** GPIO operations. */
#define LED1_ON PORTB |= (1U << LED1)                    // Enable LED1.
#define LED1_OFF PORTB &= ~(1U << LED1)                  // Disable LED1.
#define LED2_ON PORTB |= (1U << LED2)                    // Enable LED2.
#define LED2_OFF PORTB &= ~(1U << LED2)                  // Disable LED2.
#define LEDS_OFF PORTB &= ~((1U << LED1) | (1U << LED2)) // Disable both LEDs.
#define BTN1_PRESSED (PINB & (1U << BTN1))               // High if BTN1 is pressed, low otherwise.

/** Time parameters. */
#define TIMEOUT_MS 100U                            // Timeout in ms.
#define TICK_PERIOD_MS 0.064F                      // Tick period, prescaler 1024.
#define CMP_MAX (TIMEOUT_MS / TICK_PERIOD_MS - 1U) // Compare value.
#define TICK_MAX (uint16_t)(CMP_MAX + 0.5F)        // Rounded compare value.

/** LED state (true = enabled, false = disabled). */
static volatile bool led_state = false;

/** True if LED1 is lit, false otherwise. */
static volatile bool led1_lit = false;

/**
 * @brief Set up system.
 */
static void setup(void)
{
    // Configure LED1 and LED2 as outputs.
    DDRB = (1U << LED1) | (1U << LED2);

    // Configure BTN1 as input with its internal pull-up enabled.
    PORTB = (1U << BTN1);

    // Enable pin change interrupt for BTN1.
    PCICR  = (1U << PCIE0);
    PCMSK0 = (1U << BTN1);

    // Enable timer 1 in CTC mode, set timeout to 100 ms.
    TCCR1B = (1U << CS10) | (1U << CS12) | (1U << WGM12);
    OCR1A  = TICK_MAX;
    TIMSK1 = (1U << OCIE1A);

    // Enable interrupts globally.
    sei();
}

/**
 * @brief Toggle the LEDs on button press.
 */
ISR(PCINT0_vect)
{
    if (BTN1_PRESSED)
    {
        led_state = !led_state;
        if (!led_state) { LEDS_OFF; }
    }
}

/**
 * @brief Alternate which of the two LEDs is lit if enabled.
 */
ISR(TIMER1_COMPA_vect)
{
    // Check the LED state, terminate if disabled.
    if (!led_state) { return; }

    // Alternate which LED is lit.
    led1_lit = !led1_lit;

    if (led1_lit)
    {
        LED1_ON;
        LED2_OFF;
    }
    else
    {
        LED1_OFF;
        LED2_ON;
    }
}

/**
 * @brief Application entry point.
 *
 * @return 0 on termination of the program (should never occur).
 */
int main(void)
{
    setup();
    while (true) {}
    return 0;
}
\end{ccode}

\code{led_state} is \code{volatile} because it's shared: the button's ISR writes it and the timer's
ISR reads it. That follows the book's rule of making anything shared with an ISR \code{volatile},
although between two ISRs the qualifier is, strictly, the same harmless habit as on the
\code{ovf_count} of \secref{c5:a:4}: on the AVR an ISR runs to completion with interrupts
disabled, so nothing can change \code{led_state} in the middle of one, and each ISR reads it afresh
every time it runs. It becomes load-bearing the moment \code{main()}'s loop reads it too.
\code{led1_lit} is only ever touched by the timer's ISR, so it could just as well be a
\code{static} local inside it, like \code{ovf_count}.

Timer 2 mirrors Timer 0 (8-bit, Normal/CTC), with one gotcha: its \textdiv1024 prescaler bit
pattern is \code{111}, not \code{101} as on Timers 0/1.

% Section 5.2, from lectures/L05/README.md: the goals, the evaluation questions, and the next
% lecture.

\section{Review}
\label{c5:sec:review}

\needspace{6\baselineskip}
\subsection*{What you should be able to do now}
After this chapter, you should:
\begin{itemize}
  \item Explain why a busy-wait delay (\code{_delay_ms}) is unsuitable once a program needs to do
    more than one thing at a time.
  \item Be able to compute a timer's tick period from the CPU clock and a chosen prescaler.
  \item Be able to work out how many overflow interrupts are needed for a given delay, for an
    8-bit timer.
  \item Understand the practical difference between Normal Mode and CTC Mode, and when you'd reach
    for each.
  \item Know why a variable shared between an ISR and \code{main()} needs to be declared
    \code{volatile}.
\end{itemize}

\needspace{6\baselineskip}
\subsection*{Questions to test yourself}
You should be able to answer:
\begin{enumerate}
  \item At 16\,MHz with a \textdiv256 prescaler, what's the resulting tick period?
  \item What's the difference between what triggers \code{TIMERn_OVF_vect} and what triggers
    \code{TIMERn_COMPA_vect}?
  \item Why is \code{volatile} optional on \code{ovf_count} in the Timer 0 example, and when does a
    variable shared with an ISR genuinely need it?
\end{enumerate}

\needspace{6\baselineskip}
\subsection*{Looking ahead}
Chapter~\ref{c6:ch} is about serial communication:
\begin{itemize}
  \item Serial communication (UART): transmitting data to a PC instead of only driving LEDs.
  \item Combining timers with UART to report readings periodically.
\end{itemize}

% Section 5.3, from lectures/L05/appendix/b_exercises.md.

\section{Exercises}
\label{c5:sec:exercises}\label{c5:app:b}

\begin{aside}
\textbf{How to work these.} The Reflection and Calculation exercises are answered on paper, from
the chapter. The Code exercises are programs for the board: build and flash them the way
\secref{c2:app:c} (Microchip Studio) or \secref{c2:app:d} (Linux) describes, or run them in the
simulator of Appendix~\ref{sim:ch} if you have no board yet.

\textbf{How to check your work.} The Reflection and Calculation exercises, and the questions some
Code exercises ask about their design, are answered in Appendix~\ref{app:answers}. The Code
exercises have no printed solution: the board tells you whether a program does what the exercise
describes. Try each exercise before you read its answer.
\end{aside}

% ------------------------------------------------------------------------------------------
\exerciseset{Concepts}
\label{c5:set:1}

\calcexercise{Tick at \textdiv256}
\label{c5:ex:1}
At 16\,MHz with a \textdiv256 prescaler, what's the resulting tick frequency and tick period? Show
your work.

\calcexercise{Overflows for 250\,ms}
\label{c5:ex:2}
Using your answer to Exercise~\ref{c5:ex:1}, how many overflow interrupts (8-bit timer) would be
needed to time a 250\,ms delay?

% ------------------------------------------------------------------------------------------
\exerciseset{Normal vs.\ CTC}
\label{c5:set:2}

\reflexercise{When to reach for CTC}
\label{c5:ex:3}
In your own words, explain when you'd reach for CTC mode instead of Normal mode with software
overflow-counting. Is there a period that's awkward to hit precisely in Normal mode but trivial in
CTC?

% ------------------------------------------------------------------------------------------
\exerciseset{Program}
\label{c5:set:3}

\codeexercise{Blinking at 2\,Hz}
\label{c5:ex:4}
Adapt the Timer 0 blink example of \secref{c5:a:4} to blink at 2\,Hz (250\,ms on, 250\,ms off)
instead of 1\,Hz. What value changes, and what does the new \code{OVF_MAX} work out to?

\reflexercise{When \code{volatile} matters}
\label{c5:ex:5}
The chapter's Timer 0 example declares \code{ovf_count} as \code{volatile}, and \secref{c5:a:4}
says that removing the qualifier there would change nothing. Explain why. Then
suppose \code{ovf_count} moved to file scope so that \code{main()} could wait for it, in a loop
such as \code{while (10U > ovf_count) {}}. What specific failure could removing \code{volatile}
now introduce, and why might it show up in one build and not in another?

\codeexercise{A 2-second period in CTC}
\label{c5:ex:6}
Configure Timer 1 in CTC mode for a period of exactly 2 seconds at 16\,MHz with a \textdiv1024
prescaler. What value goes in \code{OCR1A}? (Timer 1 is 16-bit, so check whether your answer
fits.)

\codeexercise{Three timers, one blink}
\label{c5:ex:7}
Capstone. Two LEDs on D9/D10 blink at a speed selected by whichever of three buttons (D11--D13)
was most recently pressed, each mapped to a different timer:
\begin{itemize}
  \item D11 selects Timer 0 (1000\,ms per toggle), D12 selects Timer 1 (500\,ms per toggle), D13
    selects Timer 2 (100\,ms per toggle).
  \item Pressing a button toggles \emph{that} timer's blinking on or off. If it was off, turn it on
    and turn the other two off (only one timer drives the LEDs at a time). If it was already the
    active timer, turn it off entirely and the LEDs go dark.
  \item All three buttons are on port B and share \code{PCINT0_vect}; detect which one was just
    pressed with the \code{prev}/\code{now} edge detection from \secref{c3:a:8}.
  \item Configure all three timers to tick at the same rate (\textdiv1024 prescaler,
    $\sim$0.064\,ms per tick) so the overflow-counting logic in \secref{c5:a:4} works identically
    for each. Timer 1 needs CTC mode with \code{OCR1A} set so its compare-match period matches
    Timer 0/2's overflow period, even though its 16-bit range could hit the target directly, so all
    three interrupt handlers end up the same shape.
  \item Define a small struct to hold \code{required_interrupts} and \code{executed_interrupts} for
    one timer (see \secref{c5:a:2} and \secref{c5:a:4} for the counting pattern), and use three
    instances, one per timer, instead of three sets of loose globals.
  \item Track the currently active timer in a single variable, using either an enum
    (\code{TIMER_NONE}, \code{TIMER_0}, \code{TIMER_1}, \code{TIMER_2}) or macros with timer
    indexes.
\end{itemize}

